%% paginas/2_pre_textuais/resumo.tex
%% Elemento Pré-Textual Obrigatório (Resumo em Português) -- Conforme ABNT NBR 6028
%% Regra: Parágrafo único, espaçamento simples, justificado, de 150 a 500 palavras.
%% Palavras-chave separadas por ponto e iniciadas com maiúscula.

\begin{center}
  \textbf{\large{RESUMO}}
\end{center}

\vspace{0.8cm}

\begin{spacing}{1.0}
  \noindent
  \MakeUppercase{\meuautor}. \textbf{\meutitulo}. \meuano. \pageref{LastPage} f. \tipotrabalho\ (\grauacademicopretendido) -- \minhainstituicao, \meulocal, \meuano.
\end{spacing}

\vspace{0.8cm}

\begin{spacing}{1.0}
  \noindent
  Este trabalho tem como objetivo apresentar um modelo altamente modular e estruturado em LaTeX para a criação de documentos acadêmicos em conformidade com as mais recentes normas da Associação Brasileira de Normas Técnicas (ABNT). Através da modularização de arquivos, dividindo os componentes em pastas pré-textuais, textuais e pós-textuais, e isolando os dados de estilo em arquivos externos, torna-se muito mais fácil a manutenção, o desenvolvimento acadêmico cooperativo e o cumprimento rígido das regras de formatação. O modelo aborda a configuração precisa de margens, tipografia padrão, espaçamentos, citações longas e listas de seções. Com esta estrutura limpa e expansível, o usuário do Overleaf consegue concentrar-se exclusivamente no conteúdo de sua pesquisa, enquanto o sistema LaTeX cuida da formatação mecânica. Os resultados práticos de compilação demonstram ganho de produtividade e redução de erros comuns em trabalhos de conclusão de curso, dissertações e teses acadêmicas.
\end{spacing}

\vspace{1cm}

\noindent
\textbf{Palavras-chave:} LaTeX. ABNT. Modelo Acadêmico. Modularidade. Formatação de Trabalhos.
